\section{Algebraic data types: further details}%
\applabel{algebraic data types}

In some sense, the theoretical development in this section is a
special case of the proceedings two sections. However, it requires
some additional abstractions that make the material less accessible:
indexed families (\appref{indexed data types}), or monadic adjunctions
and distributive laws (\appref{cat dev}). Therefore, we decided to
include more specialised proofs for the simply-typed case, that use
less categorical reasoning and mostly concern sets and functions.

We start by proving that every Cartesian strength is a $\G$-strength
(cf.~\remarkref{cartesian vs symmetry strengths}).

\begin{lemma}\lemmalabel{Cartesian strengths are symmetric strengths}
  Let $\strength : X \times FY \to F(X \times Y)$ be a Cartesian
  strength for a functor $F : \Set \to \Set$. Then
  $\strength_{\carrier\G,\X} : \carrier\G \times FX \to F(\carrier\G
  \times X)$ is a $\G$-strength.
\end{lemma}
\begin{proof}
  It's more natural to use diagrammatic reasoning here:
  \diagram{15}
  \diagram{16}
  as we wanted.
\end{proof}

Next, we prove that a $\G$-strength for $F$ and an $\G$-action $\A$
let us equip $F\carrier\A$ with a $\G$-action:

\begin{proof}[of \lemmaref{strength lifts actions}]%
  This definition is a $\G$-action by straightforward calculation:
  \[
    e \act t =
    \fmap (\act) (\strength(\unit, t))
    \explain={first axiom}
    \fmap (\act) (\fmap(\fun x \pair \unit x)t)
    \explain={map fusion}
    \fmap(\fun x \unit \act x)t
    =
    \fmap(\fun x \unit\act x)t
    =
    \fmap(\id)t
    =
    t
  \]
  \begin{align*}
    g \act (h \act t) &=
    \fmap (\act) (\strength(g, \fmap (\act) (\strength(h, t))))
    \\&\explain={strength nat.}
    \fmap (\act) (\fmap (\fun{\triple{g'}{h'}{x}}\pair{g'}{h'\act x})
      (\strength(g,
    (\strength(h, t)))))
    \\&\explain={map fusion}
    \fmap (\fun{\triple{g'}{h'}x}g'\act h'\act x)
    (\strength(g,
    (\strength(h, t))))
    \\&=
    \fmap (\fun{\triple{g'}{h'}x}(g'\mul h')\act x)
    (\strength(g,
    (\strength(h, t))))
    \\&\explain={map fusion}
    \fmap (\act) (\fmap (\fun{\triple{g'}{h'}x}\pair{g'\mul h'}x)
      (\strength(g,
    (\strength(h, t)))))
    \\&\explain={second axiom}
    \fmap (\act) (\strength(g \mul h, t))
    \defines
    (g\mul h) \act t
  \end{align*}
  and so we have an action.

  Let $h : \A \equito \B$ be equivariant. Calculate:
  \begin{align*}
    \fmap h (g \act t)
    &=
    \fmap h (\fmap (\act) (\gstrength (g, t)))
    \explain={map fusion}
    \fmap(\fun{\pair{g'}x}h(g' \act x))(\gstrength (g, t))
    \\&=
    \fmap(\fun{\pair{g'}x}g' \act h(x))(\gstrength (g, t))
    \explain={map fusion}
    \fmap(\act) \fmap(\id \times h)(\gstrength (g, t))
    \\&\explain={strength nat.}
    \fmap(\act) \gstrength(g, \fmap ht)
    =
    g \act \fmap h t
  \end{align*}
  and so we have an equivariant map $\fmap h : F\A \equito F\B$.
\end{proof}

Next, we prove that $\G$-strengths for each component functor induce a
$\G$-strength for the coproduct functor:

\begin{proof}[\propositionref{modularity}]%
  Follows by $\Set$'s distributivity: the natural
  isomorphism $\carrier\G \times (\X \amalg \Y) \isomorphic
  (\carrier\G \times \X) \amalg (\carrier\G \times \Y)$.
\end{proof}

The following technical result is the workhorse that allows us to both
derive the action on the data type from its strength, and to reason
about properties of parameterised maps. It is
prevalent in Fiore et al.'s work on abstract
syntax~\cite{fiore-szamozvancev:formal-soas,%
  fiore-szamozvancev:soas-theory,%
  fiore-ranchod:substructural-ast}, originally proposed by
\citeauthor{bird-paterson:generalised-folds-for-nested-datatypes}~\cite{%
  bird-paterson:generalised-folds-for-nested-datatypes}. Here we use
the following special case.

\begin{theorem}[parameterised initiality]%
  \theoremlabel{parameterised initiality}
  Let $A$ be a set, $F : \Set \to \Set$ a functor, $\roll : F\fix F
  \to \fix F$ an initial algebra, and $\strength : A\times FX \to F(A
  \times \X)$ a natural transformation in $\X$. For every $F$-algebra
  $f : F\X \to \X$ there is a unique \emph{$A$-parameterised}
  homomorphism $\pfold f : \pair{\fix F}{\roll} \to \pair{X}{f}$, i.e., a
  function $\pfold f : A \times \fix F \to \X$ satisfying:
  \[
  \pfold f\pair a{\roll t} = f(\fmap(\strength(a,t)))
  \]
\end{theorem}
\begin{proof}%
  Using both the given strength and the Cartesian-closed structure,
  define the following morphism:
  \diagram{17}
  Currying, we have $h \definedby \curry g : F(X^A) \to X^A$, and so
  by ordinary initiality, there is a unique homomorphism into $X^A$:
  \diagram{18}
  We define the required parameterised traversal as the uncurried function:
  \[
  \pfold f \definedby \uncurry k :
  A \times \fix F \xto{\id \times k} A \times X^A \xto\eval X
  \]
  We need to show that it satisfies the required equation, and that it
  is the unique such function. Calculate:
  \diagram{19}
  Now, take any map $p : A\times \fix F \to \X$ that satisfies:
  \diagram{20}
  Let $q \definedby \curry p : \fix F \to X^A$. Calculate:
  \diagram{21}
  and so, by the universal property of the exponential, $q$ satisfies
  ($*$), and so $q = k$, hence:
  \[
  p = \uncurry q = \uncurry k = \pfold f
  \]
  and we proved uniqueness.
\end{proof}

We now lay the capstone in the theory of data types, the Equivariant
Initiality \theoremref{representation for algebras}:

\begin{proof}[of \theoremref{representation for algebras}]%
  We have three things to do:
  \begin{itemize}
    \item show the unit axiom for $(\act)$.
    \item show the associativity axiom for $(\act)$.
    \item show that $\fold f$ is equivariant when $f$ is equivariant.
  \end{itemize}
  All three are straightforward calculations, appealing to initiality
  or parameterised initiality.

  For the unit axiom, appeal to initiality:
  \diagram{23}
  and so, by initiality, $\unit \act t = t$.

  For associativity, we appeal to parameterised initiality, for the
  composite strength:
  \diagram{24}
  via the  calculation in \figref{associativity of initial algebra action}.
  \begin{figure}
    \diagram{25}
    \caption{associativity of the action over the initial algebra}
    \figlabel{associativity of initial algebra action}
  \end{figure}
  Therefore, $(g \mul h) \act t = g \act (h \mul t)$, and we have
  an action. Note that the defining property of $(\act)$ makes the
  initial algebra structure equivariant $\roll : F\fix F \equito \fix F$.

  Finally, take any equivariant algebra $f : F\A \equito \A$, and
  consider the unique $F$-homomorphism $\fold f : \fix F \to
  \carrier\A$. We appeal to parametrised initiality w.r.t.~the
  $\G$-strength, and calculate as in \figref{equivariant initiality}.
  \begin{figure}
    \diagram{26}
    \caption{equivariance of the homomorphism into an equivariant algebra}
    \figlabel{equivariant initiality}
  \end{figure}
  Therefore $\fold f(g \act t) = g \act (\fold f)$, as we wanted.
\end{proof}

The proofs from \subsecref{wlog-traversals} are identical and special
cases to their indexed variants, and do not use abstract categorical
concepts. We therefore only prove them in the indexed case
(cf.~\appref{indexed data types}).
